\section{Schemes}

\begin{exercise}\label{ex2.1}
	\textcolor{cyan}{Atiyah-Macdonald Exercise 3.21(i)}.
\end{exercise}

\begin{exercise}\label{ex2.2}
	Let $X=\bigcup_iV_i$ and each $\left(V_i,\mathcal{O}_X|_{V_i}\right)$ is an affine scheme $\operatorname{Spec}\!\left(A_i\right)$. Since $U=\bigcup_jU_{f_j}$ is the union of some basic open sets, use \Cref{ex2.1} so $\left(V_i\cap U_{f_j},\mathcal{O}_X|_{V_i\cap U_{f_j}}\right)$ is $\operatorname{Spec}\!\left(\left(A_i\right)_{f_j}\right)$.
\end{exercise}

\begin{exercise}\label{ex2.3}
	
\end{exercise}

\begin{exercise}\label{ex2.4}
	
\end{exercise}

\begin{exercise}\label{ex2.5}
	
\end{exercise}

\begin{exercise}\label{ex2.6}
	
\end{exercise}

\begin{exercise}\label{ex2.7}
	
\end{exercise}

\begin{exercise}\label{ex2.8}
	\begin{enumerate}
		\item[(ii)] Use \Cref{2.15}.
	\end{enumerate}
\end{exercise}

\begin{exercise}\label{ex2.9}
	
\end{exercise}

\begin{exercise}\label{ex2.10}
	
\end{exercise}

\begin{exercise}\label{ex2.11}
	\begin{enumerate}
		\item[(i),(ii)] Consider a residue field and pass to vector spaces. Use the right exactness of tensor product.
		\item[(iii)] If $m>n$, let $\psi:A^m\xrightarrow{\phi}A^n\hookrightarrow A^m$. By \Cref{2.4}, $\psi$ satisfies $\psi^n+a_{n-1}\psi^{n-1}+\cdots+a_0=0$. Pick such equation to have the minimum degree, then $\psi$ injective implies $a_0\neq0$. Apply the equation to element $(0,\cdots,0,1)$ to obtain contradiction.
	\end{enumerate}
\end{exercise}

\begin{exercise}\label{ex2.12}
	
\end{exercise}

\begin{exercise}\label{ex2.13}
	
\end{exercise}

\begin{exercise}\label{ex2.14}
	
\end{exercise}

\begin{exercise}\label{ex2.15}
	\begin{enumerate}[label=(\roman*),ref=\theexercise(\roman*)]\crefalias{enumi}{exercise}
		\item\label{ex2.15.1} Use \Cref{ex2.14} to increase indices.
		\item\label{ex2.15.2} \textcolor{magenta}{$x_i=\sum_k\left(y_k-\mu_{kl}\left(y_k\right)\right)$ is an equation in $\bigoplus M_i$, so $\mu_{ij}$ cannot be applied directly.} Use canonical projection $\pi_i:\bigoplus M_i\rightarrow M_i$ on the equation first.
	\end{enumerate}
\end{exercise}

\begin{exercise}[Universal property of direct limit]\label{ex2.16}
	
\end{exercise}

\begin{exercise}\label{ex2.17}
	
\end{exercise}

\begin{exercise}\label{ex2.18}
	
\end{exercise}

\begin{exercise}\label{ex2.19}
	
\end{exercise}
